% Source: https://artofproblemsolving.com/community/c3h1686730_parallelogram
% Stable identity: aops-topic-1686730-q01
Let $h$ be the distance between the parallel lines $AB$ and $CD$, and let $h_1$, $h_2$ be the distances from $P$ to $AB$ and to $CD$. Since $P$ is between the two lines, $h_1+h_2=h$.

\begin{asy}
usepackage("amsmath");
size(6.5cm);
defaultpen(fontsize(10pt));
pair A=(0,0),B=(8,0),C=(10,5),D=(2,5);
pair P=(4.2,1.6);
pair M1=(A+D)/2,M2=(B+C)/2;
draw(M1--M2,blue+dotted+linewidth(0.8));
draw(A--B--C--D--cycle,linewidth(0.9));
draw(P--A);draw(P--B);draw(P--C);draw(P--D);
pair F1=(P.x,0),F2=(P.x,5);
draw(P--F1,dashed);draw(P--F2,dashed);
real s=0.35;
draw(F1+(s,0)--F1+(s,s)--F1+(0,s),linewidth(0.6));
draw(F2+(s,0)--F2+(s,-s)--F2+(0,-s),linewidth(0.6));
dot(P);
label("$A$",A,SW);label("$B$",B,SE);label("$C$",C,NE);label("$D$",D,NW);label("$P$",P,dir(165));
label("$h_1$",(P+F1)/2,E);label("$h_2$",(P+F2)/2,E);
label("\mbox{}",(min(currentpicture,user=true)+max(currentpicture,user=true))/2,black); // keeps the text after the figure black when LaTeX typesets the labels inline
\end{asy}

Then, using $AB=CD$,
\[
[PAB]+[PCD]=\dfrac{1}{2}AB\cdot h_1+\dfrac{1}{2}CD\cdot h_2=\dfrac{1}{2}AB\cdot h=\dfrac{1}{2}[ABCD].
\]
So $[PAB]=[PCD]$ holds exactly when $h_1=h_2$, i.e.\ when $P$ lies on the line midway between $AB$ and $CD$ (the dotted line through the midpoints of $AD$ and $BC$). In the same way $[PBC]=[PDA]$ holds exactly when $P$ lies on the line midway between $BC$ and $DA$, the line through the midpoints of $AB$ and $CD$.

($\Rightarrow$) If all four areas are equal, $P$ lies on both midway lines. These two lines meet in exactly one point, the centre $O$ of the parallelogram (the midpoint of both diagonals, which lies on both midway lines because the half-turn about $O$ swaps $AB$ with $CD$ and $BC$ with $DA$). So $P=O=AC\cap BD$.

($\Leftarrow$) If $P=O$, then $P$ is on both midway lines, so $[PAB]=[PCD]$ and $[PBC]=[PDA]$. Moreover $[OAB]=\dfrac{1}{2}[ABC]=\dfrac{1}{4}[ABCD]$, because $O$ is the midpoint of $AC$, and similarly $[OBC]=\dfrac{1}{2}[ABC]=\dfrac{1}{4}[ABCD]$. Hence all four areas equal $\dfrac{1}{4}[ABCD]$. $\blacksquare$
